% ======================================================================
% Replication and Evaluation-Sensitivity Audit of
% "Pump and Dumps in the Bitcoin Era"
% Springer LNCS (llncs.cls) project for Financial Cryptography 2027
% (submission deadline 17 Sep 2026; double-blind).
%
% HOW TO BUILD (Overleaf):
%   * Upload this whole latex/ folder as a project; compiler: pdfLaTeX;
%     root document: main.tex.
%   * llncs.cls v2.24 and splncs04.bst (official Springer LNCS author
%     kit, downloaded 2026-08-01 from
%     https://resource-cms.springernature.com/springer-cms/rest/v1/content/19238648/data/)
%     are included in this folder. Overleaf also provides the class
%     built-in via its "Springer Lecture Notes in Computer Science"
%     template, so the local copy is a pinned convenience, not a hard
%     requirement. splncs04.bst is unused while references are inline
%     \bibitem entries; it is kept for a possible later switch to BibTeX.
%
% TWO BUILDS FROM ONE SOURCE:
%   * \anonymoustrue  -> anonymized FC 2027 submission build (default).
%   * \anonymousfalse -> named build (Zhennan Yu, Independent Researcher).
%   * See ANONYMIZATION_NOTES.md for the complete list of redactions.
%   * \neutralpositionalitytrue additionally neutralizes the
%     industry-experience clauses in Sect. 5.2/5.3 (default OFF: the
%     FC 2027 CFP requires no author names, affiliations,
%     acknowledgments, or obvious references, but does not clearly
%     require removing non-identifying positionality detail).
%
% Source of truth: paper/manuscript_draft_v0.5.md (2 Aug 2026).
% The draft-status banner of v0.5 is omitted as draft metadata; its
% protocol-freeze statement is retained in Sect. 8 (see notes file).
% ======================================================================
\documentclass[runningheads]{llncs}

\usepackage[T1]{fontenc}
\usepackage{graphicx}
\usepackage{amsmath}
\usepackage{amssymb}
\usepackage{booktabs}
\usepackage{url}
\usepackage{cite}          % compresses [7,8,9,10] -> [7-10]

% ----------------------- build toggles --------------------------------
\newif\ifanonymous
\anonymoustrue             % <<< FLIP HERE: \anonymoustrue = double-blind
%\anonymousfalse           %                \anonymousfalse = named build

\newif\ifneutralpositionality
\neutralpositionalityfalse % see header note and ANONYMIZATION_NOTES.md

% Repository link: redacted in the anonymized build.
\newcommand{\repolink}{%
\ifanonymous
[link redacted for review; artifacts available to reviewers on request]%
\else
\url{github.com/nathanskill/pump-and-dump-replication-audit}%
\fi}

% Sect. 5.2 industry-experience clause (removed only if
% \neutralpositionalitytrue).
\newcommand{\industryclause}{%
\ifneutralpositionality\else{} --- from several years working in the
retail-trading industry, where coordinated promotional campaigns are
observed operationally ---\fi}

% ----------------------------------------------------------------------

\begin{document}

\title{Replication and Evaluation-Sensitivity Audit of ``Pump and Dumps
in the Bitcoin Era''}
\titlerunning{Replication and Evaluation-Sensitivity Audit}

\ifanonymous
\author{Anonymous Submission to FC 2027}
\authorrunning{Anonymous Submission to FC 2027}
\institute{}
\else
\author{Zhennan Yu}
\authorrunning{Z. Yu}
\institute{Independent Researcher, Sydney, Australia\\
\email{yuzhennan323@gmail.com}}
\fi

\maketitle

\begin{abstract}
Reported scores for cryptocurrency pump-and-dump classifiers are usually
produced by row-level cross-validation, in which rows densely sampled
from the same short event window may fall on both sides of a split.
Whether this evaluation is robust to the choice of evaluation unit and
to time order has not been tested for the widely used random-forest
baseline released with La Morgia et al.'s ICCCN 2020 study. We audit
that baseline under a pre-frozen protocol and report three findings.
First, a newly written implementation reproduces the released classifier
script's aggregate precision, recall and F1 exactly at all three feature
frequencies (5\,s, 15\,s, 25\,s) on the released matrices. The published
Table~III values, however, cannot be recovered from currently public
materials: the paper-described parameters, the released code, and the
post-paper matrices differ (the same authors' journal extension tracks
the released artifact far more closely); we document the differences and
leave their cause open. Second, on an identical universe of 326 common
events, discrimination is largely insensitive to the evaluation family
across upstream-style row-level, event-exclusive, and purged
forward-time evaluation: the paired difference in average precision
between row-level and event-exclusive families has a 95\%
cluster-bootstrap interval covering zero at every frequency, and a
matched-universe comparison shows the forward family's nominally higher
score to be a test-composition effect. Third, event-level detection
nonetheless degrades at the finest frequency (event detection rate
0.64--0.92 at the released decision threshold), and ablating the four
time-of-day features lowers forward average precision by only
0.010--0.028. All conclusions are restricted to the event-centred public
benchmark; no real-market prevalence, false-alert or deployment claim is
made. Code, protocols and hash-verified aggregate artifacts are public.

\keywords{Cryptocurrency pump-and-dump \and Replication \and Evaluation
methodology \and Machine learning for security \and Fraud detection.}
% NOTE: the markdown source has no keyword list; these keywords were
% added because the LNCS format requires them -- author to confirm.
\end{abstract}

\section{Introduction}\label{sec:intro}

Pump-and-dump operations coordinate rapid purchases of thinly traded
crypto assets to attract victims into buying at inflated
prices~\cite{Kamps2018,Xu2019,Victor2019,Hamrick2021}. The market-data
detector of La Morgia et al.~\cite{LaMorgia2020} is one of the most
widely used public baselines in this space: its labelled feature
matrices and classifier script are openly released and have anchored
subsequent studies that reuse the released
artifact~\cite{Akyildirim2022,Chadalapaka2022,Moura2026,Losavio2026} as
well as adjacent detection work on independently collected
data~\cite{Nghiem2021,Hu2023}.

Reported scores for such detectors are usually produced by row-level
cross-validation over feature rows that are densely sampled from short
windows around known events. Rows from a single event share the same
episode, so a split that mixes them across training and test sets
answers a different question than a deployment setting, which concerns
\emph{new events at later times}~\cite{Kapoor2023}. To our knowledge,
whether this distinction matters \emph{for this benchmark} has not been
quantified on the released matrices and classifier. The nearest works
are the original authors' journal extension, which stress-tests its
trained detector on multi-year single-coin market streams and on later
out-of-domain crowd-pump events but contains no chronological
train-early/test-later split of the pump events~\cite{LaMorgia2023}; a
partial re-implementation that evaluates the baseline under a
non-shuffled row-level split~\cite{Chadalapaka2022}; a
chronological-split study on a rebuilt hourly panel derived from the
same upstream repository~\cite{Losavio2026}; and an unsupervised
low-latency detector whose comparison with prior work notes, in
passing, that several published detectors train on data from after the
pump so that their reported detection latencies are theoretical lower
bounds~\cite{Bello2023} --- an observation this audit systematizes.
None combines event-exclusive and purged forward-time evaluation with a
pre-frozen event-level outcome specification and cluster-bootstrap
uncertainty.

This paper asks one question: \textbf{to what extent can the released
baseline's results be traced, replicated, and shown to be robust to the
choice of evaluation family?} Contributions:

\begin{enumerate}
\item \textbf{A verified re-implementation.} A newly written
implementation whose aggregate precision, recall and F1 exactly match
the released classifier script at all three frequencies on the released
matrices (Section~\ref{sec:replication}).
\item \textbf{A documented traceability gap.} The published Table~III
values could not be recovered from the currently public artifact alone;
we record the parameter, code and matrix differences and leave their
source open (Section~\ref{sec:tableiii}).
\item \textbf{A three-family evaluation-sensitivity audit} on an
identical 326-event universe, with a pre-frozen event-level outcome
specification (eligibility, benchmark onset, interval-censored delay,
cooldown-aware alert aggregation, cluster-bootstrap uncertainty) and a
training-side calibrated threshold
(Sections~\ref{sec:threefam}--\ref{sec:ablation}).
\item \textbf{Open, hash-verified artifacts} under a conservative
redistribution boundary (Section~\ref{sec:repro}).
\end{enumerate}

This is scoped as a \emph{partial replication of the released artifact}
plus an evaluation audit in the tradition of TESSERACT and the ``dos
and don'ts'' line of work in security
ML~\cite{Rougier2017,Gundersen2018,Pineau2021,Pendlebury2019,Arp2022};
it makes no claim of being the first replication, no claim of severe
leakage or performance collapse, and no real-market deployment claim.
Following~\cite{Arp2022}, this is not a finger-pointing exercise: the
upstream artifact is one of the few public, reusable baselines in this
space, which is precisely what makes it worth auditing.

\section{Background and the Upstream Artifact}\label{sec:background}

The upstream repository~\cite{UpstreamRepo} (pinned commit
\texttt{d71250d}) releases labelled market-feature matrices at 5\,s,
15\,s and 25\,s aggregation, a feature-generation script, and a
random-forest classifier script (MIT licence). Rows carry a
\texttt{pump\_index} event key, a \texttt{date} bin label, twelve
features (order-book rush indicators, trade/volume/price statistics,
and four cyclical time-of-day encodings), and an outcome label
\texttt{gt}. Throughout, \emph{row}, upstream's \emph{chunk}, and
\emph{bin} refer interchangeably to one aggregation window of width
$\Delta f$.

Three provenance facts constrain interpretation. First, the
repository's README describes manual pump-start labelling, but the
public scripts initialize \texttt{gt=0} and do not expose the manual
labelling step; label provenance is therefore unresolved from public
materials. Second, ten event groups in the current matrices contain no
positive row; we exclude them from event-level denominators rather than
treating them as negative events. Third, the matrices cover roughly
$\pm$24 hours around known events rather than a natural market stream,
so no evaluation family on this artifact can estimate real-market
prevalence, daily false alerts, or analyst workload. An inquiry about
paper-era materials was sent to the corresponding author on 22 July
2026, with the results of this study already in hand; its outcome will
be reported here. [{\itshape STATUS AT SUBMISSION: pending --- resolves
5 August 2026 under the pre-declared 14-day rule: ``no reply at
submission time'' if unanswered}]

\section{Methods}\label{sec:methods}

\subsection{Models and Thresholds}\label{sec:models}

The primary model is the released-code random forest
(\texttt{n\_estimators=200}, \texttt{max\_depth=5},
\texttt{min\_samples\_leaf=1}, \texttt{random\_state=1}). The
paper-described configuration (\texttt{max\_depth=4},
\texttt{min\_samples\_leaf=6}) is retained as a traceability
sensitivity and never silently interchanged. Two decision thresholds
are frozen: the replication anchor $\tau = 0.5$ (the released script's
\texttt{predict()} behaviour under scikit-learn~\cite{Sklearn2011}),
and a training-calibrated $\tau^{*}$ selected only from training-side
out-of-fold probabilities by chunk-level F1 with ties resolved upward;
the outer test is evaluated once, so that threshold selection cannot
borrow information from the evaluation
side~\cite{Varma2006,Cawley2010}.

\subsection{Evaluation Families}\label{sec:evalfamilies}

All families are compared on the identical row universe of the 326
events common to all three frequencies (identical up to a
per-checkpoint purge of at most 3 events in S2; exact at the primary
checkpoint).

\begin{itemize}
\item \textbf{S0 --- row-level 5-fold} (upstream style): non-shuffled
stratified row-level cross-validation. Rows of one event may appear on
both sides of a split; S0 therefore reports chunk-level metrics only.
\item \textbf{S1 --- event-exclusive 5-fold:} \texttt{GroupKFold} by
\texttt{pump\_index}~\cite{Roberts2017}; $\tau^{*}$(fold) from
event-grouped inner out-of-fold probabilities (four inner folds); every
eligible event is scored exactly once by a model that never saw it.
\item \textbf{S2 --- purged forward-time:} common-calendar,
event-exclusive, outcome-label-free expanding assignment; events
overlapping a calendar cutoff are excluded rather than split
(\emph{purged}~\cite{LopezDePrado2018,Bergmeir2012}). $\tau^{*}$ at
each checkpoint is selected on an expanding-time inner split of the
training side (inner train fraction 0.8, with $\tau = 0.5$ as
fallback); full selection rules are recorded in the frozen run
configurations. The 80\% checkpoint is primary; the 40/50/60/70\%
checkpoints are dependent sensitivity windows and are all reported. A
disclosed pre-freeze structural pilot (Gate~0) inspected positive-event
counts, so the checkpoints are not presented as preregistered.
\end{itemize}

\subsection{Frozen Event-Level Outcomes}\label{sec:outcomes}

An event is \emph{eligible} only if its group has exactly one
\texttt{gt=1} row (317 of 327 groups per frequency; test-enforced
invariant). The one non-common group excluded per frequency in forming
the 326-event universe is itself a no-positive group, so all 317
eligible events lie inside the 326-event universe (326 = 317 eligible +
9 remaining no-positive groups). The unique positive row's
\texttt{date} is a \emph{benchmark reference onset-bin start} $t_0$,
not verified economic truth. Because \texttt{date} is a bin label of
width $\Delta f$, information is earliest available at
$t_{\mathrm{available}} = \texttt{date} + \Delta f$, so onset-bin delay
is interval-censored at $\Delta f$ and never reported as zero. Within
each held-out event, alerts (score $\geq \tau$) are emitted in time
order from the full window start under the paper's 30-minute cooldown;
an early alert may legitimately suppress an onset-adjacent one. We
report event detection rates (EDR): EDR$_{\mathrm{exact}}$ (emitted
alert sourced from the \texttt{gt=1} row); EDR@30/60/120\,s over
half-open windows $(t_0, t_0+H]$; Lead@30/60/120\,s over $[t_0-L,
t_0]$ separately; delay conditional on detection within 120\,s with
explicit denominators; and \emph{benchmark-negative alert episodes}
(emitted alerts from \texttt{gt=0} rows) with 95\% percentile
event-cluster bootstrap intervals ($n=2000$, seed 20260722; resampling
whole events to respect within-event dependence~\cite{Field2007}). The
phrases ``real false-alarm rate'', ``daily false alerts'' and ``analyst
workload'' are prohibited by protocol. Average precision --- a
precision-recall summary --- is the primary discrimination metric
throughout, because PR-based measures remain informative under the
extreme class imbalance of this benchmark where ROC-based summaries do
not~\cite{Saito2015}.

\section{Results}\label{sec:results}

\subsection{Code-Level Replication}\label{sec:replication}

The newly written implementation reproduces the released classifier
script's aggregate precision, recall and F1 exactly --- to every
reported decimal --- at all three frequencies on the released matrices
(script-hash-verified rerun; artifacts:
\texttt{gate0\_upstream\_stdout\_capture\_20260721},
\texttt{gate0\_released\_code\_rf\_all\_v2\_20260721},
\texttt{gate0\_baseline\_comparison\_20260721});
see Table~\ref{tab:replication}.

\begin{table}[t]
\centering
\caption{Code-level replication on the released matrices.
$^*$AP is computed by the re-implementation only; the released script
does not output probability scores. One row per frequency because the
two sources agree exactly.}
\label{tab:replication}
\footnotesize
\setlength{\tabcolsep}{4pt}
\begin{tabular}{llcccc}
\toprule
Frequency & Source & Precision & Recall & F1 & AP \\
\midrule
25\,s & released script / re-implementation & 0.9658 & 0.8896 & 0.9261 & 0.9644$^*$ \\
15\,s & released script / re-implementation & 0.9679 & 0.8549 & 0.9079 & 0.9472$^*$ \\
5\,s  & released script / re-implementation & 0.9433 & 0.7350 & 0.8262 & 0.9026$^*$ \\
\bottomrule
\end{tabular}
\end{table}

\subsection{Table III Traceability}\label{sec:tableiii}

We were unable to recover the published Table~III values from the
currently public artifact alone. The paper-described tree constraints,
the released code defaults, and the post-paper released matrices
differ, and the manual label-generation step is not public. The same
authors' journal extension~\cite{LaMorgia2023} reports updated
per-frequency results under its own evaluation protocol (5-fold
cross-validation; the journal states it does not use a standard
train/test split, employing a random 50/50 split only to select its
rush-order detection threshold) on a 317-event Binance collection ---
a count that exactly matches the eligible-event invariant of the
currently released matrices (Section~\ref{sec:outcomes}), which is
itself public evidence that the released matrices postdate the
conference paper. We document these differences without adjudicating
their source (artifact:
\texttt{gate0\_paper\_described\_rf\_all\_20260721}); see
Table~\ref{tab:tableiii}.

\begin{table}[t]
\centering
\caption{Published Table~III values against the currently public
artifact.}
\label{tab:tableiii}
\footnotesize
\setlength{\tabcolsep}{4pt}
\begin{tabular}{llccc}
\toprule
Frequency & Config & Precision & Recall & F1 \\
\midrule
25\,s & published Table III & 0.937 & 0.913 & 0.918 \\
25\,s & released artifact (script defaults) & 0.9658 & 0.8896 & 0.9261 \\
25\,s & paper-described trees on current matrices & 0.9685 & 0.8738 & 0.9187 \\
15\,s & published Table III & 0.913 & 0.844 & 0.877 \\
15\,s & released artifact (script defaults) & 0.9679 & 0.8549 & 0.9079 \\
15\,s & paper-described trees on current matrices & 0.9713 & 0.8549 & 0.9094 \\
5\,s  & published Table III & 0.924 & 0.784 & 0.840 \\
5\,s  & released artifact (script defaults) & 0.9433 & 0.7350 & 0.8262 \\
5\,s  & paper-described trees on current matrices & 0.9604 & 0.6877 & 0.8015 \\
25\,s & journal version, own protocol (5-fold CV)~\cite{LaMorgia2023} & 0.982 & 0.912 & 0.945 \\
15\,s & journal version, own protocol (5-fold CV)~\cite{LaMorgia2023} & 0.964 & 0.849 & 0.900 \\
5\,s  & journal version, own protocol (5-fold CV)~\cite{LaMorgia2023} & 0.946 & 0.729 & 0.824 \\
\bottomrule
\end{tabular}
\end{table}

The deviations run in both directions (the released artifact exceeds
the published F1 at 25\,s and 15\,s but trails it at 5\,s). The journal
rows~\cite{LaMorgia2023} track the released-artifact rows far more
closely than Table~III does, consistent with a benign artifact-drift
reading in which the conference paper reports an earlier, smaller
matrix generation; the split protocols nonetheless differ, so the
correspondence is indicative rather than a replication target, and the
drift's source remains undetermined from public materials.

\subsection{Three Evaluation Families, Identical
Universe}\label{sec:threefam}

\begin{table}[t]
\centering
\caption{Three evaluation families on the identical universe. Cells
give AP [95\% cluster CI]; S0/S1 pooled OOF, S2 single forward test.
S0 here is refit on the 326-event common universe; its AP therefore
differs in the third decimal from the full-matrix values of
Section~\ref{sec:replication}.}
\label{tab:families}
\scriptsize
\setlength{\tabcolsep}{3.5pt}
\begin{tabular}{lccc}
\toprule
Family & 25\,s & 15\,s & 5\,s \\
\midrule
S0 row-level 5-fold & 0.9649 [0.949, 0.979] & 0.9462 [0.925, 0.965] & 0.9025 [0.876, 0.929] \\
S1 event-exclusive 5-fold & 0.9636 [0.947, 0.978] & 0.9466 [0.925, 0.966] & 0.9007 [0.873, 0.927] \\
S2 forward @80\% (primary) & 0.9781 [0.948, 0.999] & 0.9572 [0.917, 0.991] & 0.9072 [0.845, 0.957] \\
\bottomrule
\end{tabular}
\end{table}

Discrimination is largely insensitive to the evaluation family
(Table~\ref{tab:families}): event-exclusive and forward-time evaluation
do not collapse the row-level result on this benchmark
(Fig.~\ref{fig:pr}: the S0 and S1 precision--recall curves are nearly
coincident at every frequency, and the S2 forward curve, on its own
test subset, tracks them). Under a shared event resample, the paired
$\Delta$AP(S1 $-$ S0) is $-0.0013$ $[-0.0039, +0.0009]$ at 25\,s,
$+0.0004$ $[-0.0022, +0.0030]$ at 15\,s and $-0.0017$ $[-0.0051,
+0.0012]$ at 5\,s --- every interval covers zero. The forward family's
nominally higher AP is a test-composition effect, not a family effect:
restricting S1's out-of-fold predictions to exactly the 66 events of
the S2 primary test set yields AP 0.9785 / 0.9591 / 0.9148 versus S2's
0.9781 / 0.9572 / 0.9072 (analysis addendum A4; artifact
\texttt{a4\_addendum\_v1\_20260723}). Cross-frequency AP differences
are not interpreted as model-performance differences (test prevalences
differ: 7.8, 6.6 and $4.9 \times 10^{-4}$ at the S2 primary
checkpoint); the same caution applies across families, whose test
universes differ --- the matched-universe comparison above addresses it
directly.

Event-level detection at the anchor threshold (EDR@120\,s): S1 pooled
0.886 / 0.845 / 0.710; S2 primary 0.922 / 0.859 / 0.641 (25\,s /
15\,s / 5\,s). At the S2 primary checkpoint EDR@120 coincides with
EDR$_{\mathrm{exact}}$ because every within-120\,s detection
originates at the onset bin itself (Section~\ref{sec:checkpoints}:
median delay equals one bin width). (Artifacts:
\texttt{formal\_s1\_s0\_v1\_20260722},
\texttt{formal\_forward\_v1\_20260722}.)

\begin{figure}[t]
\centering
\includegraphics[width=\textwidth]{figures/fig1_pr_curves}
\caption{Precision--recall by evaluation family and feature frequency.
Three panels (25\,s / 15\,s / 5\,s). In each, the S0 row-level and S1
event-exclusive out-of-fold curves (solid) are computed on the full
326-event universe; the S2 forward @80\% curve (dashed) is on its
66-event terminal test subset. S0 and S1 are nearly coincident
throughout; discrimination does not collapse under event-exclusive or
forward-time evaluation. (Generated by \texttt{src/make\_figures.py}
from the per-row prediction artifacts, which are not redistributed; see
the Data Availability statement.)}
\label{fig:pr}
\end{figure}

\subsection{Forward-Time Checkpoints (S2, All
Reported)}\label{sec:checkpoints}

\begin{table}[t]
\centering
\caption{Forward-time checkpoints (S2, all reported). Cell format:
AP / $\tau^{*}$ / EDR@120\,s. The 80\% checkpoint is primary.}
\label{tab:checkpoints}
\footnotesize
\setlength{\tabcolsep}{3pt}
\begin{tabular}{lccc}
\toprule
Checkpoint & 25\,s (AP / $\tau^{*}$ / EDR@120) & 15\,s (AP / $\tau^{*}$ / EDR@120) & 5\,s (AP / $\tau^{*}$ / EDR@120) \\
\midrule
40\% & .9479 / .287 / .821 & .9400 / .311 / .842 & .8585 / .395 / .663 \\
50\% & .9413 / .312 / .866 & .9250 / .375 / .822 & .8461 / .354 / .637 \\
60\% & .9508 / .132 / .898 & .9340 / .321 / .835 & .8702 / .163 / .654 \\
70\% & .9712 / .471 / .914 & .9637 / .475 / .839 & .9232 / .306 / .710 \\
\textbf{80\% (primary)} & \textbf{.9781 / .704 / .922} & \textbf{.9572 / .506 / .859} & \textbf{.9072 / .338 / .641} \\
\bottomrule
\end{tabular}
\end{table}

At the primary checkpoint (Table~\ref{tab:checkpoints}),
EDR$_{\mathrm{exact}}$ at the anchor threshold is 0.922 [95\% cluster
CI 0.844--0.984] (25\,s), 0.859 [0.766--0.938] (15\,s) and 0.641
[0.516--0.750] (5\,s). The calibrated threshold cuts both ways: at
5\,s, $\tau^{*} = 0.338$ raises EDR$_{\mathrm{exact}}$ to 0.734 while
benchmark-negative episodes rise from 0.11 to 0.14 per event, but at
25\,s, $\tau^{*} = 0.704$ \emph{lowers} EDR$_{\mathrm{exact}}$ from
0.922 to 0.875. $\tau^{*}$ itself is unstable across checkpoints
(25\,s: 0.132 at 60\% versus 0.704 at 80\%) --- an expected symptom of
score-scale shift between the inner model that selects the threshold
and the larger outer model it is applied to. Median detection delay
equals one bin width at every frequency --- the interval-censored
minimum: detections occur at the onset bin itself.
Checkpoint-to-checkpoint variation is non-monotone --- the 70\%
checkpoint exceeds the primary at 15\,s and 5\,s --- and all
checkpoints are reported (Fig.~\ref{fig:checkpoints}, left). The
$\tau^{*}$ instability is visible in Fig.~\ref{fig:checkpoints}
(right): at 25\,s it swings from 0.13 at the 60\% checkpoint to 0.70 at
80\%. (Artifacts: \texttt{formal\_forward\_v1\_20260722}.)

Per protocol, the full frozen outcome set is reported in analysis
addendum A4 (artifact \texttt{a4\_addendum\_v1\_20260723}): EDR@30 =
EDR@60 = EDR@120 in every family $\times$ frequency $\times$ threshold
cell; Lead@120 ranges from 0.000 (25\,s forward) to 0.125 (5\,s forward
at $\tau^{*}$); delay median, IQR and p90 all equal one bin width;
detection denominators $n_{\mathrm{detected}}/n_{\mathrm{eligible}}$
are 59/64, 55/64 and 41/64 at the forward primary anchor (281/317,
268/317 and 225/317 for S1 pooled, 25\,s/15\,s/5\,s); and
benchmark-negative episode means carry 95\% cluster intervals (5\,s
forward: 0.11 [0.05, 0.19] at the anchor, 0.14 [0.06, 0.23] at
$\tau^{*}$).

\begin{figure}[t]
\centering
\includegraphics[width=\textwidth]{figures/fig2_checkpoints}
\caption{Forward-time checkpoints: detection and threshold stability.
Left: event detection rate @120\,s at the anchor threshold across the
five forward checkpoints; the 5\,s frequency is uniformly lowest and
the variation is non-monotone. Right: the calibrated threshold
$\tau^{*}$ by checkpoint, illustrating its instability (e.g.\ 25\,s,
0.13 at 60\% to 0.70 at 80\%) --- an expected symptom of
inner-to-outer model score-scale shift. (Generated by
\texttt{src/make\_figures.py} from the per-row prediction artifacts,
which are not redistributed; see the Data Availability statement.)}
\label{fig:checkpoints}
\end{figure}

\subsection{Time-of-Day Feature Ablation (S2)}\label{sec:ablation}

Removing the four cyclical time-of-day columns and re-selecting
$\tau^{*}$ on the ablated model's own inner validation
(Table~\ref{tab:ablation}):

\begin{table}[t]
\centering
\caption{Time-of-day feature ablation at the 80\% checkpoint (S2).}
\label{tab:ablation}
\footnotesize
\setlength{\tabcolsep}{3pt}
\begin{tabular}{lccc}
\toprule
Checkpoint (80\%) & 25\,s & 15\,s & 5\,s \\
\midrule
AP full $\rightarrow$ ablated & .9781 $\rightarrow$ .9612 & .9572 $\rightarrow$ .9470 & .9072 $\rightarrow$ .8805 \\
EDR@120 (anchor) full $\rightarrow$ ablated & .922 $\rightarrow$ .875 & .859 $\rightarrow$ .797 & .641 $\rightarrow$ .641 \\
\bottomrule
\end{tabular}
\end{table}

Across all fifteen combinations the AP reduction lies between 0.010 and
0.028: the time-of-day features contribute a small, consistent
increment and are not a dominant shortcut. By protocol, this is a
descriptive sensitivity --- a drop does not prove leakage and a null
effect does not prove its absence. The full fifteen-combination grid is
given in Appendix~\ref{app:ablation} (artifact:
\texttt{formal\_forward\_ablation\_v1\_20260722}).

\section{Discussion}\label{sec:discussion}

\subsection{Reading the Agreement between Evaluation
Families}\label{sec:agreement}

The three families are best read as three different measurement
instruments pointed at the same released artifact, not as three
attempts at one deployment quantity. Row-level cross-validation asks
whether pump bins are separable when the model may have seen rows from
the same event; event-exclusive validation asks whether the separation
transfers to unseen events; the forward families add calendar order. On
this benchmark the instruments largely agree
(Section~\ref{sec:threefam}). The honest reading is narrow: within
these $\pm$24-hour event windows, the order-book and volume features
that distinguish onset bins from surrounding background are strong
enough that event memorization is not the main driver of the reported
scores. Had the forward families collapsed --- an outcome this protocol
was equally prepared to report --- the row-level numbers would have had
to be read as instrument artifacts. They did not collapse here.
Agreement between instruments inside the benchmark, however, does not
certify the benchmark: none of the families can speak for a natural
market stream.

\subsection{The 5-Second Degradation}\label{sec:fivesecond}

Mechanically, the finest grid is the hardest one: at 5\,s the test
prevalence is lowest, each bin aggregates the fewest trades, per-bin
statistics are noisiest, and the pump impulse is spread over more bins,
so the single onset bin is less reliably distinctive ---
EDR$_{\mathrm{exact}}$ falls to 0.641 at the anchor threshold, and the
calibrated threshold buys back part of it (0.734) at a small
benchmark-negative cost. Beyond the mechanics, the author's
reading\industryclause{} is that a pump is a burst produced by a
coordinated group, not a smooth signal, and a 5-second grid is neither
intrinsically fine nor coarse: whether the burst lands sharply inside
one bin depends on the coordination behind it --- the group's size, its
scripting, how quickly its participants act in concert --- and that
structure varies across events and is invisible in exchange market
data. This reading is interpretive, not established by the artifact.

\subsection{What These Results Mean --- and Do Not Mean --- for
Practitioners}\label{sec:practitioners}

Two boundaries matter more than any score. First, a market-data
detector sees the footprint, not the actor. The accounts and projects
involved are effectively anonymous, and the interest networks behind
promoted tokens are complex and opaque; nothing in these metrics
identifies who is behind an event, and nothing here supports
account-level accusation. Second,
\ifneutralpositionality
in the author's assessment of promotion channels,
\else
in the author's experience of promotion channels,
\fi
much of the visible ``community'' around such tokens is itself part of
the operation: coordinated accounts posting near-identical promotional
messages and staged momentum. The social layer around a pump is
manipulated in the same way the order book is. A benchmark-negative
episode rate on the order of one episode per ten events on these
curated windows must therefore never be read as a real-world
false-alert rate --- a natural stream contains organic volatility and
staged chatter that this artifact cannot represent, and the true
operational rate is neither measured nor bounded here. Results of this
kind support triage: deciding where limited human attention should go,
under review, with the burden of proof left where it belongs.

\section{Limitations}\label{sec:limitations}

Conclusions are bounded by the event-centred universe: $\pm$24-hour
windows around known events cannot estimate real-market prevalence,
daily false alerts, deployment calibration or analyst workload. Label
provenance is unresolved from public materials; \texttt{gt=0} is never
interpreted as confirmed absence of a pump, and the ten no-positive
groups stay outside event denominators. The published-table gap is
documented; its cause is left open. Results concern one dataset, one
model family and one market period; the disclosed Gate-0 pilot means
the forward checkpoints are not label-unseen.

\section{Conclusion}\label{sec:conclusion}

This audit asked whether the results of a widely reused pump-and-dump
baseline survive scrutiny of their evaluation unit, and the answer is
largely yes, with one telling exception. A newly written implementation
reproduces the released classifier script exactly at all three
frequencies, although the published Table~III values remain
unrecoverable from currently public materials --- a traceability gap we
document without adjudicating. On an identical 326-event universe,
discrimination is largely insensitive to the evaluation family: the
paired difference in average precision between row-level and
event-exclusive evaluation has a 95\% cluster-bootstrap interval
covering zero at every frequency, and the forward family's nominally
higher score is a test-composition effect, not a family effect.
Event-level detection nonetheless degrades at the finest frequency ---
EDR$_{\mathrm{exact}}$ at the anchor threshold falls from 0.922 at
25\,s to 0.641 at 5\,s at the primary forward checkpoint --- a weakness
that row-level scores alone do not reveal. All of this holds only
inside the event-centred benchmark; no real-market prevalence,
false-alert or deployment claim follows. Studies that reuse this
artifact should report event-exclusive and event-level metrics
alongside row-level scores: on this benchmark the added cost of doing
so is near zero, and the information gained at 5\,s is substantial.

\section{Reproducibility Statement}\label{sec:repro}

All protocols, code, tests, run configurations, split manifests,
aggregate metrics and artifact hashes are public at \repolink; each
results subsection above names its artifact directory. The outcome
specification was frozen before the outer runs at the commit now
identified as \texttt{c2736ed}. Archived run configurations record the
same commit under its pre-rewrite identifier (\texttt{23089ce}): on 23
July 2026, after the runs completed, the repository's commit
\emph{metadata} (author identity; assistant trailers) was rewritten;
amendment A3 records the authoritative old-to-new SHA map and the
tree-hash-equality verification showing that file contents are
unchanged commit-for-commit. Amendments A1/A2 changed claim boundaries
only. Formal runs used Python 3.9 with scikit-learn 1.6.1, numpy
1.26.4, pandas 2.2.3 (pinned in \texttt{requirements.txt}). Upstream
matrices, raw transactions, the Telegram group list and the event file
are not redistributed; reproducers obtain them from the original
repository at the pinned commit and verify against recorded SHA-256
hashes. The novelty statement in Section~\ref{sec:intro} rests on a
citation-forward screen of works citing~\cite{LaMorgia2020} conducted
on 21 July 2026, together with the literature search for this draft.

\paragraph{Data availability.} See above: derived aggregate artifacts
(metrics, split manifests, selection tables) and artifact hashes are
public; upstream inputs are obtained from~\cite{UpstreamRepo} at the
pinned commit. Per-row prediction files --- the inputs to
Fig.~\ref{fig:pr} and to the A4 cluster bootstrap --- are not
redistributed: they are regenerated deterministically (fixed seeds)
from the pinned upstream inputs by \texttt{src/run\_formal\_cv.py} and
\texttt{src/run\_formal\_forward.py}, and their SHA-256 hashes are
recorded in
\texttt{artifacts/a4\_addendum\_v1\_20260723/artifact\_manifest.csv}
for verification. A results-complete release (tag
\texttt{v0.2.0-results-freeze}, whose tree contains the frozen
protocol, all artifact directories cited in
Sections~\ref{sec:results}--\ref{sec:discussion}, the A4 addendum, and
\texttt{CITATION.cff}) is additionally archived under a persistent DOI
(Zenodo), providing a tamper-evident anchor independent of the
rewritten Git history and of the upstream repository's continued
availability.
% NOTE: do NOT insert a real Zenodo DOI in the anonymized build -- the
% Zenodo record names the author (see ANONYMIZATION_NOTES.md).
[{\itshape Zenodo DOI to be inserted at preprint.}]

\paragraph{Ethics.} This study analyses only the publicly released
research artifact~\cite{UpstreamRepo}; no new data were collected from
any platform or individual. Nothing in the analysis identifies accounts
or individuals or supports account-level attribution, and the upstream
Telegram-derived materials are not redistributed.

\paragraph{AI-assistance disclosure.} Large-language-model assistants
(Anthropic Claude) were used under the author's direction for code
drafting, literature search, data processing, and manuscript editing,
including drafting the Discussion from the author's dictated analysis.
All protocol decisions, frozen specifications, claims, and the
Discussion's analysis are the author's; all reported numbers were
verified against hash-recorded artifacts. [{\itshape Adapt to venue
disclosure format at submission.}]

% ---------------------------------------------------------------------
\begin{thebibliography}{28}

\bibitem{LaMorgia2020}
M. La Morgia, A. Mei, F. Sassi, J. Stefa. \emph{Pump and Dumps in the
Bitcoin Era: Real Time Detection of Cryptocurrency Market
Manipulations.} ICCCN 2020. DOI 10.1109/ICCCN49398.2020.9209660.

\bibitem{UpstreamRepo}
SystemsLab-Sapienza, \emph{pump-and-dump-dataset}, GitHub repository,
commit \texttt{d71250d} (MIT licence).

\bibitem{LaMorgia2023}
M. La Morgia, A. Mei, F. Sassi, J. Stefa. \emph{The Doge of Wall
Street: Analysis and Detection of Pump and Dump Cryptocurrency
Manipulations.} ACM Transactions on Internet Technology
23(1):11:1--11:28, 2023. DOI 10.1145/3561300 (arXiv:2105.00733).

\bibitem{Sklearn2011}
F. Pedregosa et al. \emph{Scikit-learn: Machine Learning in Python.}
JMLR 12:2825--2830, 2011.

\bibitem{LopezDePrado2018}
M. L\'opez de Prado. \emph{Advances in Financial Machine Learning.}
Wiley, 2018.

\bibitem{Rougier2017}
N. P. Rougier et al. \emph{Sustainable computational science: the
ReScience initiative.} PeerJ Computer Science 3:e142, 2017.

\bibitem{Kamps2018}
J. Kamps, B. Kleinberg. \emph{To the moon: defining and detecting
cryptocurrency pump-and-dumps.} Crime Science 7:18, 2018. DOI
10.1186/s40163-018-0093-5.

\bibitem{Xu2019}
J. Xu, B. Livshits. \emph{The Anatomy of a Cryptocurrency
Pump-and-Dump Scheme.} USENIX Security 2019, 1609--1625.
arXiv:1811.10109.

\bibitem{Victor2019}
F. Victor, T. Hagemann. \emph{Cryptocurrency Pump and Dump Schemes:
Quantification and Detection.} IEEE ICDMW 2019. DOI
10.1109/ICDMW.2019.00045.

\bibitem{Hamrick2021}
J. T. Hamrick et al. \emph{An examination of the cryptocurrency
pump-and-dump ecosystem.} Information Processing \& Management
58(4):102506, 2021. DOI 10.1016/j.ipm.2021.102506.

\bibitem{Nghiem2021}
H. Nghiem, G. Muric, F. Morstatter, E. Ferrara. \emph{Detecting
cryptocurrency pump-and-dump frauds using market and social signals.}
Expert Systems with Applications 182:115284, 2021. DOI
10.1016/j.eswa.2021.115284.

\bibitem{Akyildirim2022}
E. Akyildirim, M. Gambara, J. Teichmann, S. Zhou. \emph{Applications of
Signature Methods to Market Anomaly Detection.} arXiv:2201.02441,
2022.

\bibitem{Chadalapaka2022}
V. Chadalapaka, K. Chang, G. Mahajan, A. Vasil. \emph{Crypto Pump and
Dump Detection via Deep Learning Techniques.} arXiv:2205.04646, 2022.

\bibitem{Hu2023}
S. Hu, Z. Zhang, S. Lu, B. He, et al. \emph{Sequence-Based Target Coin
Prediction for Cryptocurrency Pump-and-Dump.} Proceedings of the ACM on
Management of Data 1(1):6:1--6:19, 2023. DOI 10.1145/3588686.

\bibitem{Moura2026}
M. S. Moura, L. Baroni, E. Ogasawara, D. S. Mendon\c{c}a. \emph{HDP+:
Leveraging Anomaly and Change Point Detection for Pump-and-Dump in
Cryptocurrency Exchanges.} Journal of Information and Data Management
17(1):194--203, 2026. DOI 10.5753/jidm.2026.5728.

\bibitem{Losavio2026}
L. Losavio, L. Persia, M. Sathe, D. Pasadakis. \emph{Fraud Detection in
Cryptocurrency Markets with Spatio-Temporal Graph Neural Networks.}
arXiv:2604.24590, 2026.

\bibitem{Varma2006}
S. Varma, R. Simon. \emph{Bias in error estimation when using
cross-validation for model selection.} BMC Bioinformatics 7:91, 2006.
DOI 10.1186/1471-2105-7-91.

\bibitem{Field2007}
C. A. Field, A. H. Welsh. \emph{Bootstrapping clustered data.} Journal
of the Royal Statistical Society: Series B 69(3):369--390, 2007. DOI
10.1111/j.1467-9868.2007.00593.x.

\bibitem{Cawley2010}
G. C. Cawley, N. L. C. Talbot. \emph{On Over-fitting in Model Selection
and Subsequent Selection Bias in Performance Evaluation.} JMLR
11:2079--2107, 2010.

\bibitem{Bergmeir2012}
C. Bergmeir, J. M. Ben\'itez. \emph{On the use of cross-validation for
time series predictor evaluation.} Information Sciences 191:192--213,
2012. DOI 10.1016/j.ins.2011.12.028.

\bibitem{Roberts2017}
D. R. Roberts et al. \emph{Cross-validation strategies for data with
temporal, spatial, hierarchical, or phylogenetic structure.} Ecography
40(8):913--929, 2017. DOI 10.1111/ecog.02881.

\bibitem{Gundersen2018}
O. E. Gundersen, S. Kjensmo. \emph{State of the Art: Reproducibility in
Artificial Intelligence.} AAAI 2018. DOI 10.1609/aaai.v32i1.11503.

\bibitem{Pineau2021}
J. Pineau et al. \emph{Improving Reproducibility in Machine Learning
Research (A Report from the NeurIPS 2019 Reproducibility Program).}
JMLR 22(164):1--20, 2021. arXiv:2003.12206.

\bibitem{Kapoor2023}
S. Kapoor, A. Narayanan. \emph{Leakage and the reproducibility crisis
in machine-learning-based science.} Patterns 4(9):100804, 2023. DOI
10.1016/j.patter.2023.100804.

\bibitem{Pendlebury2019}
F. Pendlebury, F. Pierazzi, R. Jordaney, J. Kinder, L. Cavallaro.
\emph{TESSERACT: Eliminating Experimental Bias in Malware
Classification across Space and Time.} USENIX Security 2019, 729--746.

\bibitem{Arp2022}
D. Arp, E. Quiring, F. Pendlebury, A. Warnecke, F. Pierazzi, C.
Wressnegger, et al. \emph{Dos and Don'ts of Machine Learning in
Computer Security.} USENIX Security 2022, 3971--3988.

\bibitem{Saito2015}
T. Saito, M. Rehmsmeier. \emph{The Precision-Recall Plot Is More
Informative than the ROC Plot When Evaluating Binary Classifiers on
Imbalanced Datasets.} PLOS ONE 10(3):e0118432, 2015. DOI
10.1371/journal.pone.0118432.

\bibitem{Bello2023}
A. S. Bello, J. Schneider, R. Di Pietro. \emph{LLD: A Low Latency
Detection Solution to Thwart Cryptocurrency Pump \& Dumps.} IEEE ICBC
2023. DOI 10.1109/ICBC56567.2023.10174922.

\end{thebibliography}

% ---------------------------------------------------------------------
\appendix

\section{Full Time-of-Day Ablation Grid (S2, All Fifteen
Combinations)}\label{app:ablation}

Removing the four cyclical time-of-day columns and re-selecting
$\tau^{*}$ on the ablated model's own inner validation; artifact:
\texttt{formal\_forward\_ablation\_v1\_20260722} (full model:
\texttt{formal\_forward\_v1\_20260722}). See
Table~\ref{tab:ablationgrid}.

\begin{table}[t]
\centering
\caption{Full time-of-day ablation grid (S2, all fifteen
combinations).}
\label{tab:ablationgrid}
\scriptsize
\setlength{\tabcolsep}{2.5pt}
\begin{tabular}{llccccccc}
\toprule
Freq. & Ckpt. & AP (full) & AP (ablated) & $\Delta$AP &
\begin{tabular}[c]{@{}c@{}}EDR@120 anchor\\(full)\end{tabular} &
\begin{tabular}[c]{@{}c@{}}EDR@120 anchor\\(ablated)\end{tabular} &
$\tau^{*}$ (full) & $\tau^{*}$ (ablated) \\
\midrule
25\,s & 40\% & 0.9479 & 0.9322 & $-0.0157$ & 0.821 & 0.805 & 0.287 & 0.178 \\
25\,s & 50\% & 0.9413 & 0.9204 & $-0.0209$ & 0.866 & 0.803 & 0.312 & 0.125 \\
25\,s & 60\% & 0.9508 & 0.9336 & $-0.0171$ & 0.898 & 0.850 & 0.132 & 0.131 \\
25\,s & 70\% & 0.9712 & 0.9575 & $-0.0137$ & 0.914 & 0.871 & 0.471 & 0.254 \\
25\,s & 80\% & 0.9781 & 0.9612 & $-0.0169$ & 0.922 & 0.875 & 0.704 & 0.582 \\
15\,s & 40\% & 0.9400 & 0.9211 & $-0.0189$ & 0.842 & 0.800 & 0.311 & 0.303 \\
15\,s & 50\% & 0.9250 & 0.9096 & $-0.0153$ & 0.822 & 0.796 & 0.375 & 0.292 \\
15\,s & 60\% & 0.9340 & 0.9163 & $-0.0177$ & 0.835 & 0.811 & 0.321 & 0.288 \\
15\,s & 70\% & 0.9637 & 0.9502 & $-0.0136$ & 0.839 & 0.817 & 0.475 & 0.426 \\
15\,s & 80\% & 0.9572 & 0.9470 & $-0.0103$ & 0.859 & 0.797 & 0.506 & 0.458 \\
5\,s  & 40\% & 0.8585 & 0.8375 & $-0.0210$ & 0.663 & 0.637 & 0.395 & 0.366 \\
5\,s  & 50\% & 0.8461 & 0.8241 & $-0.0220$ & 0.637 & 0.618 & 0.354 & 0.325 \\
5\,s  & 60\% & 0.8702 & 0.8446 & $-0.0255$ & 0.654 & 0.630 & 0.163 & 0.280 \\
5\,s  & 70\% & 0.9232 & 0.8955 & $-0.0277$ & 0.710 & 0.667 & 0.306 & 0.260 \\
5\,s  & 80\% & 0.9072 & 0.8805 & $-0.0267$ & 0.641 & 0.641 & 0.338 & 0.288 \\
\bottomrule
\end{tabular}
\end{table}

\end{document}
